\documentclass[11pt]{article}
\usepackage{amsmath,amssymb,geometry}
\geometry{margin=25mm}
\begin{document}
\section*{Core model: two-stage aircraft recovery with precancellation}
Source: Section 3.1, PDF pp.4--5, Table 3 and (1)--(9).
The complete route-based BM is retained; the nonanticipative first-stage
precancellation and all scenario recovery decisions form one model.
\paragraph{Data and feasible routes.}
$L,A,S,B$ are flights, aircraft, disruption scenarios and required end bases.
$p_s>0$, $\sum_s p_s=1$. Costs $c_l$ and $d_l^s>c_l$ are advance and
real-time cancellation costs. $c_b\ge0$ penalizes base imbalance and
$n_b$ is the required number at base $b$. $O_s$ indexes airport/interval
slot restrictions with integer capacity $u_o^s$.
For aircraft $a$ and scenario $s$, $R_a^s$ is the \emph{complete} set of
physically feasible routes on its finite time-space network, with every route satisfying the network requirements below. Each route
starts at the aircraft's actual available airport/time, uses compatible
flight copies (scheduled time plus an allowed nonnegative delay), and
connects successive flights only when airports agree and the previous
arrival plus required turn time does not exceed the next departure.
It respects aircraft availability/disruption restrictions, terminates
within the recovery horizon, and contains at most one copy of each flight.
Ground waiting arcs connect consecutive events at the same airport.
An idle route can represent an aircraft staying at its starting base.
These rules define $R_a^s$ from the supplied flight times, delay-copy grid,
aircraft availability, compatibility and turn-time data.

For each route, $\delta_{lr}$ indicates inclusion of flight $l$,
$\theta_{or}$ counts its departures/arrivals in slot $o$, and
$\gamma_{br}$ indicates its end base. Its cost is
\[
 q_{ar}=\sum_{l\in L}\delta_{lr}(k_{al}+\varphi_{lr}),
\]
where $k_{al}$ is aircraft swapping cost and $\varphi_{lr}$ the delay
cost of the selected copy. These are route coefficients.
\paragraph{Variables.}
$x_l\in\{0,1\}$ is advance cancellation and is common to every scenario;
$y_l^s\in\{0,1\}$ is real-time cancellation;
$z_{ar}^s\in\{0,1\}$ selects a route; $v_b^s\in\mathbb Z_+$ is the
nonnegative base-shortage slack. The source prints $\mathbb Z$ in (9),
but defines $v$ as a penalized slack for missing aircraft. This document
makes that physical nonnegative domain explicit.
\paragraph{Single complete model.}
\begin{align}
\min\quad &\sum_{l\in L}c_lx_l+
 \sum_{s\in S}p_s\left(\sum_{l\in L}d_l^s y_l^s+
 \sum_{a\in A}\sum_{r\in R_a^s}q_{ar}z_{ar}^s+
 \sum_{b\in B}c_bv_b^s\right),\tag{1}\\
\text{s.t.}\quad &\sum_{a\in A}\sum_{r\in R_a^s}\delta_{lr}z_{ar}^s+x_l+y_l^s=1
 &&l\in L,s\in S,\tag{2}\\
&\sum_{r\in R_a^s}z_{ar}^s\le1 &&a\in A,s\in S,\tag{3}\\
&\sum_{a\in A}\sum_{r\in R_a^s}\theta_{or}z_{ar}^s\le u_o^s
 &&o\in O_s,s\in S,\tag{4}\\
&\sum_{a\in A}\sum_{r\in R_a^s}\gamma_{br}z_{ar}^s+v_b^s\ge n_b
 &&b\in B,s\in S,\tag{5}\\
&x_l,y_l^s,z_{ar}^s\in\{0,1\},\qquad v_b^s\in\mathbb Z_+.
\end{align}

\end{document}
